\section{实验结果与分析}
\mbox{}\par
\subsection{实验一：证明有效性与效率}
\subsubsection{证明完成率}
为评价不同引理选择方法的证明有效性，本文分别在包含 78 个任务的主实验组和包含 31 个任务的困难实验组上，比较 Prefix Search（PS）、Classical Houdini（CH）和 LLM-GLS 的证明完成情况。各方法使用相同的任务、固定候选引理池、形式化后端与时间预算。只有目标属性及所选引理的初始性和归纳性均得到形式化证明、最终返回 \texttt{PROVED} 的任务计入成功；\texttt{TIMEOUT}、\texttt{UNKNOWN} 和 \texttt{ERROR} 分别统计，并保留在任务总数中。方法 $M$ 在实验组 $G$ 上的证明完成率定义为
\[
  \mathrm{CR}(M,G)=
  \frac{N_{\mathrm{PROVED}}(M,G)}{N_G}\times 100\%,
\]
其中 $N_G$ 为该组任务数。实验一对各方法独立重复运行三轮，证明完成率与后续证明成本分析均使用这同一批运行记录。完成率按轮对全部任务计算，再报告三轮的均值与标准差，避免将重复运行视为新增任务。

\begin{table}[t]
  \centering
  \caption{主实验组与困难实验组的证明完成情况}
  \label{tab:proof_completion_rate}
  \small
  \setlength{\tabcolsep}{4pt}
  \begin{tabular}{lrrrrr}
    \hline
    方法 & 成功 & 超时 & 未决 & 错误 & 完成率（\%） \\
    \hline
    \multicolumn{6}{l}{主实验组（78 个任务）} \\
    PS      & 62 & 0 & 14 & 2 & 79.49 $\pm$ 0.00 \\
    CH      & 66 & 0 & 10 & 2 & 84.62 $\pm$ 0.00 \\
    LLM-GLS & 71 & 0 & 5  & 2 & 91.03 $\pm$ 0.00 \\
    \hline
    \multicolumn{6}{l}{困难实验组（31 个任务）} \\
    PS      & 20 & 2 & 3 & 6 & 64.52 $\pm$ 0.00 \\
    CH      & 20 & 2 & 3 & 6 & 64.52 $\pm$ 0.00 \\
    LLM-GLS & 23 & 1 & 1 & 6 & 74.19 $\pm$ 0.00 \\
    \hline
  \end{tabular}
  \par\smallskip
  \begin{minipage}{\columnwidth}
    \footnotesize
    注：成功、超时、未决、错误分别对应 \texttt{PROVED}、\texttt{TIMEOUT}、\texttt{UNKNOWN}、\texttt{ERROR}。状态数量为三轮运行的均值，完成率以三轮均值 $\pm$ 标准差表示。
  \end{minipage}
\end{table}

表~\ref{tab:proof_completion_rate} 汇总了两组实验结果。三轮运行的完成率标准差均为 $0.00$ 个百分点。在主实验组中，PS、CH 和 LLM-GLS 分别完成了 62、66 和 71 个任务，证明完成率分别为 79.49\%、84.62\% 和 91.03\%。LLM-GLS 相较 PS 和 CH 分别提高了 11.54 和 6.41 个百分点。这表明在主实验组和给定预算下，LLM-GLS 的完整流程比两种基线完成了更多证明任务。

在困难实验组中，PS 和 CH 均完成了 20 个任务，完成率均为 64.52\%；LLM-GLS 完成了 23 个任务，完成率为 74.19\%，比两种基线均多完成 3 个任务，按未四舍五入的完成率计算提高 9.68 个百分点。与主实验组相比，三种方法在困难实验组的完成率均有所下降，但 LLM-GLS 在该组仍保持最高完成率。这一结果表明，LLM-GLS 的完成率优势在预先划定的困难实验组中也得到观察。

按三轮运行的均值统计未完成任务的终止状态，主实验组中，PS、CH 和 LLM-GLS 分别有 14、10 和 5 个任务返回 \texttt{UNKNOWN}，各有 2 个任务返回 \texttt{ERROR}，均无 \texttt{TIMEOUT}。困难实验组中，PS 和 CH 各有 2 个 \texttt{TIMEOUT}、3 个 \texttt{UNKNOWN} 和 6 个 \texttt{ERROR}；LLM-GLS 有 1 个 \texttt{TIMEOUT}、1 个 \texttt{UNKNOWN} 和 6 个 \texttt{ERROR}。由于同组各方法的 \texttt{ERROR} 数量相同，LLM-GLS 在主实验组增加的证明成功数对应更少的 \texttt{UNKNOWN}，在困难实验组则对应更少的 \texttt{TIMEOUT} 和 \texttt{UNKNOWN}。这些状态区分了预算不足、证明未决和工具执行异常，不能统一解释为目标属性不成立或候选池中不存在有效引理组合。因此，上述完成率反映的是完整流程在给定预算下的证明覆盖范围。下一小节在三种方法共同成功的任务上比较证明时间与形式化查询量，以避免不同成功任务构成影响成本比较。

\subsubsection{证明时间与形式化查询量}
在三种方法均完成证明的 78 个任务中，主实验组有 58 个，困难实验组有 20 个。证明时间包含选择、形式化检查及模型推理等实际开销，候选生成时间作为共用前置成本单独处理，不作为对比的指标。表~\ref{tab:proof_time_queries} 汇总了各方法在这些任务对应的 234 条记录上的均值。LLM-GLS 的平均证明时间为 9.70 s，相较 PS 的 37.88 s 和 CH 的 11.23 s，分别减少 74.39\% 和 13.61\%。

\begin{table}[t]
  \centering
  \caption{共同成功任务上的证明成本比较}
  \label{tab:proof_time_queries}
  \small
  \setlength{\tabcolsep}{4pt}
  \begin{tabular}{lrrr}
    \hline
    方法 & 证明时间（s） & 查询次数 & LLM 调用次数\\
    \hline
    PS      & 37.88 & 20.46 & 0\\
    CH      & 11.23 & 7.72 & 0\\
    LLM-GLS & 9.70  & 5.68 & 1.13\\
    \hline
  \end{tabular}
  \par\smallskip
  \begin{minipage}{\columnwidth}
    \footnotesize
    注：各项为三轮运行记录的均值；每种方法包含相同的 78 个任务、共 234 条记录。
  \end{minipage}
\end{table}

在同一任务集合上，LLM-GLS 的平均形式化查询次数为 5.68 次，相较 PS 的 20.46 次和 CH 的 7.72 次，分别减少 72.24\% 和 26.41\%。在这些共同成功任务上，LLM-GLS 平均调用 1.13 次 LLM。

由于证明时间已经包含模型推理，表~\ref{tab:proof_time_queries} 中较短的平均证明时间表明，新增模型调用没有抵消这一任务集合上的总体时间收益。

综合表1的证明完成率，以及表2在共同成功任务上的平均证明时间和形式化查询次数，实验一表明：在固定候选池与统一预算下，LLM-GLS 在主实验组和困难实验组均完成更多证明任务；

在三种方法共同成功的任务上，其平均证明时间更短、形式化查询更少。非前缀组合选择与结构化证明反馈各自的贡献，留待后续关键设计对照实验分析。

\subsection{实验二：候选池规模与干扰鲁棒性}
\mbox{}\par
\subsubsection{候选池规模的影响}
\mbox{}\par
\subsubsection{干扰候选类型的影响}
\mbox{}\par
\subsection{实验三：关键设计对照}
\mbox{}\par
\subsubsection{推荐组合与排序前缀对照}
\mbox{}\par
\subsubsection{结构化证明反馈对照}
\mbox{}\par
\subsubsection{修复轮次与调用成本}
\mbox{}\par
\subsection{典型任务与未决案例分析}
\mbox{}\par
